\documentclass[11pt]{article}
\usepackage[margin=1in]{geometry}
\usepackage{amsmath, amssymb, amsthm, mathtools}
\usepackage{enumitem}
\usepackage{booktabs}
\usepackage{array}
\usepackage{hyperref}

\theoremstyle{plain}
\newtheorem{theorem}{Theorem}[section]
\newtheorem{proposition}[theorem]{Proposition}
\newtheorem{lemma}[theorem]{Lemma}
\newtheorem{corollary}[theorem]{Corollary}
\newtheorem{conjecture}[theorem]{Conjecture}
\theoremstyle{definition}
\newtheorem{definition}[theorem]{Definition}
\newtheorem{remark}[theorem]{Remark}

\newcommand{\cA}{\mathcal{A}}
\newcommand{\cB}{\mathcal{B}}
\newcommand{\cZ}{\mathcal{Z}}
\newcommand{\cV}{\mathcal{V}}
\newcommand{\cN}{\mathcal{N}}
\newcommand{\C}{\mathbb{C}}
\newcommand{\Z}{\mathbb{Z}}
\DeclareMathOperator{\rank}{rank}

\title{Locality from Hidden Zeros:\\
Cyclic $1$-zeros uniquely determine $\phi^3$ tree amplitudes}
\author{Jonathan Valenzuela\\
\small Department of Engineering Physics, Tsinghua University}
\date{\today}

\begin{document}
\maketitle

\begin{abstract}
We prove that the cyclic $1$-zero conditions of Rodina uniquely determine the
ordered $\phi^3$ tree amplitude $\cA_n^{\mathrm{tree}}$ within the full space of
rational functions with simple poles in planar Mandelstam variables---a space that
includes all possible non-local (crossing-chord) contributions. At $n=5$ and $n=6$
the proof is carried out by explicit symbolic computation: every non-local term is
eliminated by the $1$-zeros, and the tree amplitude emerges as the unique survivor.
At $n=7$ the result is confirmed numerically. In all cases, \emph{locality}
(poles only in planar channels) is a \emph{derived consequence} of the $1$-zero
conditions, not an assumption. This is in contrast to Rodina's original proof,
which works within the space of Feynman-diagram sums and therefore assumes locality
from the outset.
\end{abstract}

\tableofcontents

%======================================================================
\section{Introduction}
%======================================================================

Rodina \cite{Rodina2024} proved that ordered tree-level amplitudes in $\mathrm{Tr}(\phi^3)$
theory are uniquely characterised by their \emph{hidden zeros}: the amplitude
$\cA_n^{\mathrm{tree}}$ vanishes on each of the $n$ cyclic $1$-zero loci
$\cZ_1, \ldots, \cZ_n$ in the space of planar Mandelstam variables. Rodina's proof
works within the space of local (Feynman-diagram) ansatze and establishes that the
$1$-zeros fix all relative coefficients---i.e., they prove \emph{unitarity} given
locality.

We ask the stronger question: do the $1$-zeros also enforce \emph{locality} itself?
That is, if we start with the broadest possible rational ansatz---including terms with
poles in non-planar (crossing-chord) channels and even higher-order poles---do the
$1$-zeros force all non-local terms to vanish?

We prove that the answer is yes at $n = 5$ and $n = 6$ by explicit computation, and
verify it numerically at $n = 7$. The proofs are purely linear-algebraic: the $1$-zero
conditions impose a system of linear equations on the ansatz coefficients, and we show
that the nullspace of this system is exactly $1$-dimensional, spanned by the tree
amplitude.

\subsection{Main results}

\begin{theorem}[Informal]
For $n = 5, 6, 7$, let $B$ be a rational function of the planar Mandelstam variables,
homogeneous of mass dimension $-(n-3)$, with at most simple poles. If $B$ vanishes on
each of the $n$ cyclic $1$-zero loci, then $B$ is proportional to the $\phi^3$ tree
amplitude:
\[
B \;=\; c \cdot \cA_n^{\mathrm{tree}}.
\]
In particular, every non-planar pole is eliminated by the $1$-zero conditions.
\end{theorem}

The formal statements are Theorems~\ref{thm:n5} and \ref{thm:n6} below.

%======================================================================
\section{Preliminaries}
\label{sec:prelim}
%======================================================================

\subsection{Planar Mandelstam variables}

Let $p_1, \ldots, p_n$ be massless on-shell momenta satisfying $\sum p_i = 0$ and
$p_i^2 = 0$. Define $X_{ij} = (p_i + p_{i+1} + \cdots + p_{j-1})^2$ for
$1 \le i < j \le n$, with $X_{ij} = X_{ji}$, $X_{ii} = 0$, and
$X_{i,i+1} = X_{1n} = 0$. At $n$ points there are $d = n(n-3)/2$ independent planar
variables forming the kinematic space $\cV_n$.

\subsection{The $1$-zero loci}

For each $k \in \{1, \ldots, n\}$, the \emph{cyclic $1$-zero at leg $k$} is the
codimension-$(n-3)$ linear subspace $\cZ_k \subset \cV_n$ cut out by
\begin{equation}
c_{k, k+2} = c_{k, k+3} = \cdots = c_{k, k+n-2} = 0,
\quad\text{where}\quad
c_{ij} := X_{ij} + X_{i+1,j+1} - X_{i,j+1} - X_{i+1,j}.
\label{eq:cdef}
\end{equation}

\subsection{The ansatz space}

\begin{definition}
The \emph{ansatz space} $\cB_n$ is the vector space of rational functions of the
$d = n(n-3)/2$ planar Mandelstam variables that are homogeneous of degree $-(n-3)$
and have at most simple poles in each variable independently. A basis for $\cB_n$ is
\[
\Big\{\, \tfrac{1}{X_{a_1} \cdots X_{a_{n-3}}} \;:\;
a_1 < \cdots < a_{n-3},\; a_i \in \{X_{ij}\}_{1 \le i < j \le n} \,\Big\},
\]
so $\dim \cB_n = \binom{d}{n-3}$.
\end{definition}

At $n=5$: $d = 5$, $\dim \cB_5 = \binom{5}{2} = 10$.
At $n=6$: $d = 9$, $\dim \cB_6 = \binom{9}{3} = 84$.
At $n=7$: $d = 14$, $\dim \cB_7 = \binom{14}{4} = 1001$.

\begin{remark}
At $n = 5$ we also consider a broader ansatz $V_5^{(-2)}$ of dimension $15$
that includes double poles $1/X_{ij}^2$. The theorem holds for both ansatze.
\end{remark}

\subsection{The tree amplitude}

The $\phi^3$ tree amplitude $\cA_n^{\mathrm{tree}} \in \cB_n$ is the sum over all
$C_{n-2}$ (Catalan number) non-crossing triangulations of the convex $n$-gon:
\[
\cA_n^{\mathrm{tree}} = \sum_{\text{triangulations}} \frac{1}{X_{a_1} \cdots X_{a_{n-3}}}.
\]
Rodina's hidden-zero theorem states that $\cA_n^{\mathrm{tree}}$ vanishes on every
$\cZ_k$.

%======================================================================
\section{Part I: The $n=5$ base case}
\label{sec:n5}
%======================================================================

\subsection{Setup}

At $n=5$ the planar variables are $X_{13}, X_{14}, X_{24}, X_{25}, X_{35}$. The five
$1$-zero loci are listed in Table~\ref{tab:loci5}.

\begin{table}[ht]
\centering
\small
\begin{tabular}{c|l|l}
\toprule
$k$ & Substitution & Free variables \\
\midrule
$1$ & $X_{13} = -X_{25},\; X_{14} = X_{24} - X_{25}$ & $X_{24}, X_{25}, X_{35}$ \\
$2$ & $X_{13} = X_{35}-X_{25},\; X_{24} = X_{25} - X_{35}$ & $X_{14}, X_{25}, X_{35}$ \\
$3$ & $X_{13} = X_{14} + X_{35},\; X_{24} = -X_{35}$ & $X_{14}, X_{25}, X_{35}$ \\
$4$ & $X_{14} = -X_{35},\; X_{24} = X_{25} - X_{35}$ & $X_{13}, X_{25}, X_{35}$ \\
$5$ & $X_{13} = X_{35}-X_{25},\; X_{14} = -X_{25}$ & $X_{24}, X_{25}, X_{35}$ \\
\bottomrule
\end{tabular}
\caption{The five cyclic $1$-zero loci at $n=5$.}
\label{tab:loci5}
\end{table}

We work with the $15$-dimensional ansatz $V_5^{(-2)}$ (including double poles), with
basis $\{1/(X_{ij} X_{kl})\}_{(ij) \ne (kl)} \cup \{1/X_{ij}^2\}$ and coefficients
$c_1, \ldots, c_{15}$ as in Table~\ref{tab:basis5}.

\begin{table}[ht]
\centering
\small
\begin{tabular}{clcccl}
\toprule
& term & & & & term \\
\midrule
$c_1$  & $1/(X_{13} X_{14})$ & & & $c_9$    & $1/(X_{24} X_{35})$ \\
$c_2$  & $1/(X_{13} X_{24})$ & & & $c_{10}$ & $1/(X_{25} X_{35})$ \\
$c_3$  & $1/(X_{13} X_{25})$ & & & $c_{11}$ & $1/X_{13}^2$ \\
$c_4$  & $1/(X_{13} X_{35})$ & & & $c_{12}$ & $1/X_{14}^2$ \\
$c_5$  & $1/(X_{14} X_{24})$ & & & $c_{13}$ & $1/X_{24}^2$ \\
$c_6$  & $1/(X_{14} X_{25})$ & & & $c_{14}$ & $1/X_{25}^2$ \\
$c_7$  & $1/(X_{14} X_{35})$ & & & $c_{15}$ & $1/X_{35}^2$ \\
$c_8$  & $1/(X_{24} X_{25})$ & & & & \\
\bottomrule
\end{tabular}
\caption{The $15$-dimensional ansatz at $n=5$.}
\label{tab:basis5}
\end{table}

\subsection{Leg-by-leg analysis}

\paragraph{Leg 1.} Substitute $X_{13} \to -X_{25}$, $X_{14} \to X_{24} - X_{25}$,
clear denominators. Row reduction yields $9$ independent equations:
\begin{equation}
\begin{aligned}
c_1 - c_5 - c_6 &= 0, & c_4 - c_{10} &= 0, \\
c_2 + c_5 - c_8 &= 0, & c_7 &= 0, \\
c_3 - c_{11} - c_{14} &= 0, & c_9 &= 0, \\
c_{12} &= 0, & c_{13} &= 0, \\
c_{15} &= 0. &&
\end{aligned}
\label{eq:leg1-5}
\end{equation}

After leg $1$: $6$-dimensional survivor space.

\paragraph{Leg 2.} Four new independent equations:
\begin{equation}
c_5 - c_{10} = 0, \qquad c_6 = 0, \qquad c_8 - c_{10} - c_{11} = 0, \qquad c_{14} = 0.
\label{eq:leg2-5}
\end{equation}

After leg $2$: $2$-dimensional survivor space, spanned by
\begin{equation}
B\big|_{\text{after legs 1,2}}
= c_{10}\, \cA_5^{\mathrm{tree}}
+ c_{11} \cdot \cN,
\quad\text{where}\quad
\cN = \frac{(X_{13} + X_{24})(X_{13} + X_{25})}{X_{13}^2\, X_{24}\, X_{25}}.
\label{eq:2d-survivor}
\end{equation}

The non-local survivor $\cN$ contains crossing-chord products $1/(X_{13} X_{24})$,
$1/(X_{13} X_{25})$ and the squared propagator $1/X_{13}^2$.

\paragraph{Leg 3.} Since $\cA_5^{\mathrm{tree}}$ vanishes on $\cZ_3$,
\[
B\big|_{\cZ_3} = -c_{11} \cdot \frac{X_{14}(X_{14} + X_{25} + X_{35})}{X_{25}\, X_{35}\, (X_{14}+X_{35})^2} = 0
\quad \Longrightarrow \quad c_{11} = 0.
\]

\paragraph{Legs 4, 5.} Redundant: $B = c_{10}\, \cA_5^{\mathrm{tree}}$ already vanishes
on $\cZ_4$ and $\cZ_5$.

\subsection{The theorem at $n=5$}

\begin{theorem}\label{thm:n5}
Let $B \in V_5^{(-2)}$. If $B$ vanishes on all five cyclic $1$-zero loci, then
$B = c \cdot \cA_5^{\mathrm{tree}}$.
\end{theorem}

\begin{proof}
The system \eqref{eq:leg1-5}--\eqref{eq:leg2-5} plus $c_{11} = 0$ comprises
$14$ independent equations in $15$ unknowns. Nullspace dimension $= 1$, spanned by
$\cA_5^{\mathrm{tree}}$.
\end{proof}

\begin{corollary}
Exactly $3$ of the $5$ cyclic $1$-zeros are independent as constraints on
$V_5^{(-2)}$.
\end{corollary}

\begin{remark}[Where non-local terms die]
Crossing-chord terms $c_7 = 1/(X_{14}X_{35})$ and $c_9 = 1/(X_{24}X_{35})$ die at
leg~$1$. The crossing-chord terms $c_2, c_3$ and the squared propagator $c_{11}$
survive two legs but die at leg~$3$. After leg~$3$, the surviving coefficients
$c_1 = c_4 = c_5 = c_8 = c_{10}$ are exactly the $5$ non-crossing triangulations
of the pentagon, all with equal coefficient.
\end{remark}

%======================================================================
\section{Part II: The $n=6$ case}
\label{sec:n6}
%======================================================================

\subsection{Setup}

At $n=6$ the $9$ planar variables are
$X_{13}, X_{14}, X_{15}, X_{24}, X_{25}, X_{26}, X_{35}, X_{36}, X_{46}$.
The simple-pole ansatz $\cB_6$ has $\binom{9}{3} = 84$ basis elements. The $6$
cyclic $1$-zero loci are listed in Table~\ref{tab:loci6}.

\begin{table}[ht]
\centering
\small
\begin{tabular}{c|l|l}
\toprule
$k$ & Substitution & Free variables \\
\midrule
$1$ & $X_{13}=-X_{26},\; X_{14}=X_{24}-X_{26},\; X_{15}=X_{25}-X_{26}$
    & $X_{24},X_{25},X_{26},X_{35},X_{36},X_{46}$ \\
$2$ & $X_{13}=X_{36}-X_{26},\; X_{24}=X_{26}-X_{36},\; X_{25}=X_{26}+X_{35}-X_{36}$
    & $X_{14},X_{15},X_{26},X_{35},X_{36},X_{46}$ \\
$3$ & $X_{13}=X_{14}+X_{36}-X_{46},\; X_{24}=X_{46}-X_{36},\; X_{35}=X_{36}-X_{46}$
    & $X_{14},X_{15},X_{25},X_{26},X_{36},X_{46}$ \\
$4$ & $X_{14}=X_{15}+X_{46},\; X_{24}=X_{25}+X_{46},\; X_{35}=-X_{46}$
    & $X_{13},X_{15},X_{25},X_{26},X_{36},X_{46}$ \\
$5$ & $X_{15}=-X_{46},\; X_{25}=X_{26}-X_{46},\; X_{35}=X_{36}-X_{46}$
    & $X_{13},X_{14},X_{24},X_{26},X_{36},X_{46}$ \\
$6$ & $X_{13}=X_{36}-X_{26},\; X_{14}=X_{46}-X_{26},\; X_{15}=-X_{26}$
    & $X_{24},X_{25},X_{26},X_{35},X_{36},X_{46}$ \\
\bottomrule
\end{tabular}
\caption{The six cyclic $1$-zero loci at $n=6$.}
\label{tab:loci6}
\end{table}

\subsection{Leg-by-leg rank analysis}

\begin{theorem}\label{thm:n6}
The cumulative constraint rank evolves as follows:
\begin{center}
\begin{tabular}{lcccc}
\toprule
Legs imposed & Raw eqs & Rank (of $84$) & Nullspace & New rank \\
\midrule
Leg $1$          & $63$  & $51$ & $33$ & $51$ \\
Legs $1, 2$      & $154$ & $76$ & $8$  & $25$ \\
Legs $1, 2, 3$   & $245$ & $81$ & $3$  & $5$  \\
Legs $1,2,3,4$   & $308$ & $82$ & $2$  & $1$  \\
Legs $1,\ldots,5$& $371$ & $83$ & $1$  & $1$  \\
All $6$          & $434$ & $83$ & $1$  & $0$  \\
\bottomrule
\end{tabular}
\end{center}
The nullspace is $1$-dimensional, spanned by $\cA_6^{\mathrm{tree}}$ (the sum over
all $C_4 = 14$ non-crossing triangulations of the hexagon, with equal coefficients).
\end{theorem}

\begin{proof}
Verified by exact symbolic computation in \texttt{sympy}. At each step, the $\cZ_k$
substitution is applied to the $84$-term ansatz, denominators are cleared, and
monomial coefficients in the $6$ free variables are set to zero.
The nullspace vector has exactly $14$ nonzero entries (all equal), matching the $14$
triangulations.
\end{proof}

\begin{corollary}
At $n=6$, exactly $5$ of the $6$ cyclic $1$-zeros are independent. Leg~$6$ is
redundant. The number of independent legs is $n - 1$ at both $n = 5$ and $n = 6$.
\end{corollary}

\subsection{Structural features}

\begin{remark}[Free-variable grading]
On each $\cZ_k$, the $6$ free variables split into \emph{fully free}
($3$ variables not appearing in any constraint expression) and
\emph{constrained-free} ($3$ variables appearing in the substitution right-hand
sides). Since the fully-free variables are genuinely independent parameters,
any rational function vanishing on $\cZ_k$ can be graded by powers of these
variables, and each grade must vanish independently. This is the mechanism behind
leg~$1$'s large rank contribution ($51$ of $84$).
\end{remark}

\begin{remark}[Intermediate structure]
After $3$ legs, the nullspace is $3$-dimensional. Unlike at $n=5$ (where the
$2$-dimensional intermediate cleanly separated into tree $+$ non-local), the three
basis vectors at $n=6$ each \emph{mix} tree and non-tree terms. Legs~$4$ and $5$
contribute one new equation each, ``unmixing'' these three vectors to isolate the
tree amplitude. This structural difference means a general-$n$ proof cannot rely on
clean tree/non-tree separation at each inductive step.
\end{remark}

%======================================================================
\section{The $n=7$ case (numerical)}
\label{sec:n7}
%======================================================================

\begin{proposition}
At $n = 7$, the ansatz space $\cB_7$ has $\binom{14}{4} = 1001$ basis elements.
Numerical computation (SVD with multiple random seeds) gives constraint rank $= 1000$,
nullspace dimension $= 1$, with the unique nullspace vector proportional to
$\cA_7^{\mathrm{tree}}$ (the sum over all $C_5 = 42$ triangulations of the heptagon).
The singular-value gap between the $1000$-th and $1001$-st singular values exceeds
$13$ orders of magnitude.
\end{proposition}

%======================================================================
\section{Discussion}
\label{sec:discussion}
%======================================================================

\subsection{Summary of results}

\begin{center}
\begin{tabular}{cccccc}
\toprule
$n$ & $\dim \cB_n$ & Rank & Nullspace & Method & Locality derived? \\
\midrule
$5$ & $15$   & $14$   & $1$ & Symbolic & Yes \\
$6$ & $84$   & $83$   & $1$ & Symbolic & Yes \\
$7$ & $1001$ & $1000$ & $1$ & Numerical & Yes \\
\bottomrule
\end{tabular}
\end{center}

\subsection{Locality is derived}

The central contribution of this work is showing that \emph{locality is a consequence
of the $1$-zero conditions}, not an independent assumption. Rodina's original proof
works within the space of Feynman-diagram sums (local ansatze) and proves that the
$1$-zeros fix all relative coefficients---establishing unitarity given locality. Our
ansatz $\cB_n$ is strictly larger: it includes every rational function with the
correct mass dimension and pole structure, including crossing-chord products that
cannot arise from any Feynman diagram. The $1$-zeros eliminate all such non-local
terms.

\subsection{The conjecture}

\begin{conjecture}\label{conj:main}
For all $n \ge 5$,
\[
\dim \ker\!\big(\text{all $n$ cyclic $1$-zeros on } \cB_n\big) = 1,
\]
with the unique generator being $\cA_n^{\mathrm{tree}}$.
\end{conjecture}

Three consecutive base cases ($n = 5, 6, 7$) are verified. Several structural features
support the conjecture:
\begin{itemize}[nosep]
\item The free-variable grading mechanism is uniform in $n$: at each leg, the fully-free
      variables (those not entering any constraint expression) form an increasingly
      large independent set, and the grading kills an increasing fraction of the
      ansatz. At $n = 5$ one leg kills $60\%$; at $n = 8$ one leg kills $\sim 98\%$.
\item The number of independent legs appears to be $n - 1$ at each $n$ tested.
\item The ratio (total conditions from all loci) / (needed conditions) grows
      superexponentially with $n$, so the system is massively overdetermined at large
      $n$.
\end{itemize}

A proof of Conjecture~\ref{conj:main} would constitute a derivation of both locality
and unitarity from the hidden-zero conditions alone.

\subsection{Relation to Rodina's proof}

Rodina's proof \cite{Rodina2024} uses the $D$-subset decomposition to organize
Feynman-diagram terms by boundary-propagator structure under a BCFW shift, then proves
by induction that the $1$-zeros fix all coefficients within the local ansatz. The key
difference with the present work:

\begin{center}
\begin{tabular}{lcc}
\toprule
& Rodina & This work \\
\midrule
Ansatz space & Local (Feynman diagrams) & Full $\cB_n$ (all poles) \\
Locality & Assumed & Derived \\
Unitarity & Derived & Derived \\
Method & Induction via $D$-subsets & Direct linear algebra \\
Scope & All $n$ & $n = 5, 6, 7$ (conjectured for all $n$) \\
\bottomrule
\end{tabular}
\end{center}

The two approaches are complementary: Rodina's result holds for all $n$ but assumes
locality; ours derives locality but (so far) only at specific $n$.

\subsection{Toward a general proof}

A proof of Conjecture~\ref{conj:main} likely requires one of:

\begin{enumerate}[nosep]
\item \textbf{Representation theory.} The cyclic group $\Z_n$ acts on both $\cB_n$ and
the constraint matrix. Decomposing by $\Z_n$-representations may reduce the problem.
\item \textbf{Commutative algebra.} The intersection of ideals
$I(\cZ_1) \cap \cdots \cap I(\cZ_n)$ at the appropriate degree controls the
nullspace. Hilbert function methods on arrangements of linear subspaces may apply.
\item \textbf{Combinatorics of the associahedron.} The tree amplitude is the
canonical form of the kinematic associahedron. The $1$-zero loci are specific
faces/sections of this polytope, and the uniqueness may follow from the facial
structure.
\end{enumerate}

\begin{thebibliography}{9}
\bibitem{Rodina2024}
L.~Rodina, \emph{Hidden zeros are equivalent to enhanced ultraviolet scaling and
lead to unique amplitudes in $\mathrm{Tr}(\phi^3)$ theory}, arXiv:2406.04234v5 (2024).
\end{thebibliography}

\end{document}
